% Versioned draft v4, 29 August 2026.
% Initial-submission candidate prepared for Discussiones Mathematicae Graph
% Theory (DMGT). The frozen, externally compiled v3 source is unchanged.
% Author metadata and declarations were filled on 29 August 2026. This file
% still requires external human review and a real compile/log/page audit.
\documentclass[11pt,a4paper]{article}
\usepackage[T1]{fontenc}
\usepackage[utf8]{inputenc}
\usepackage{lmodern}
\usepackage{amsmath,amssymb,amsthm}
\usepackage[mathlines]{lineno}
\usepackage[margin=1in]{geometry}
\usepackage{longtable,booktabs,array}
\IfFileExists{microtype.sty}{\usepackage{microtype}}{}
\IfFileExists{xurl.sty}{\usepackage{xurl}}{\usepackage{url}}
\IfFileExists{needspace.sty}{\usepackage{needspace}}{}
\providecommand{\Needspace}[1]{}
\usepackage[hidelinks]{hyperref}
\hypersetup{
  pdftitle={Dual General Position in the Join of a Complete Graph and a Tree},
  pdfauthor={Yi Yuteng},
  hidelinks,
  pdfcreator={LaTeX}}
\urlstyle{same}
\setlength{\emergencystretch}{3em} % prevent overfull lines
\providecommand{\tightlist}{%
  \setlength{\itemsep}{0pt}\setlength{\parskip}{0pt}}
\newcolumntype{L}[1]{>{\raggedright\arraybackslash}p{#1}}
\setcounter{secnumdepth}{3}
\newtheorem{theorem}{Theorem}[section]
\newtheorem{lemma}[theorem]{Lemma}
\newtheorem{corollary}[theorem]{Corollary}
\newtheorem{proposition}[theorem]{Proposition}
\numberwithin{equation}{section}

\title{Dual General Position in the Join of a Complete Graph and a Tree}
\author{%
  Yi Yuteng%
  \thanks{Corresponding author. Email:
  \href{mailto:yiyuteng29@163.com}{\texttt{yiyuteng29@163.com}}.}\\
  \small Independent Researcher\\
  \small Shanghai, China}
\date{}

\begin{document}
\maketitle
\linenumbers

\begin{abstract}

Let \(r\ge 1\) and let \(T\) be a tree of order at least three. We
classify the dual general-position sets of the mixed join \(K_r+T\)
according to whether they meet the universal clique. The apex-meeting
branch is governed by partitions of \(T\) into two induced cliques,
while the apex-avoiding branch is governed by a tree parameter
\(\beta(T)\), where
\(\beta(T)=\max\{|X|:X\subseteq V(T),\ \Delta(T[X])\le 1,
|N_T(x)\setminus X|\le 1\text{ for every }x\in X\}\). This gives

\[
\operatorname{gp}_d(K_r+T)=
\begin{cases}
r+2,&T\in\{P_3,P_4\},\\
\beta(T),&\text{otherwise}.
\end{cases}
\]

We characterize every \(\beta\)-feasible set by the degrees and labels
of neighboring vertices, and give a linear-time, linear-space
rooted-tree dynamic program that reconstructs a maximum set. Worked
families include stars, once-subdivided stars, and paths. Two bounded
computations with distinct checking logic test the implementation against
exhaustive local search and a definition-first shortest-path checker; these
calculations are supporting evidence, not proof.
\end{abstract}

\noindent\textbf{Keywords:} dual general position; graph join; tree;
graph convexity; dynamic programming.

\medskip
\noindent\textbf{2020 Mathematics Subject Classification:}
05C12 (primary); 05C05, 05C69, 05C76, 05C85 (secondary).

\section{Introduction}\label{sec:introduction}

A set of vertices of a connected graph is in \emph{dual general
position} if it is in general position and its complement is convex.
Equivalently, no geodesic contains three selected vertices, and every
geodesic between two unselected vertices stays outside the selected set.
Tian and Klav\v{z}ar introduced and developed this variant of the
general-position problem and proved the criterion that combines these
two conditions~\cite{tian-klavzar-variety}. For a recent overview of
general-position problems and their variants, see~\cite{chandran-et-al-survey}.

The standard general-position problem has also been studied under graph
operations. In particular, Ghorbani et al.~\cite{ghorbani-et-al-operations}
determined the standard general-position number for joins and several
other operations. Those results concern standard general position and do
not determine the dual general-position number considered here.

The behavior of dual general position under graph products was
subsequently studied by Dokyeesun, Klav\v{z}ar, Kuziak, and Tian~\cite{dokyeesun-et-al-products}. In
particular, their work posed the problem of determining the dual general
position number of lexicographic products with a complete first factor.
Jiang's 2026 preprint~\cite[Theorem 3.2]{jiang-complete-first-factor} gives the formula

\[
\operatorname{gp}_d(K_m\circ G)=m q_2(G), \qquad m\ge 2,
\]

for every nonempty finite simple graph \(G\), where \(q_2(G)\) is the
largest size of one class in a partition of \(V(G)\) into two sets that
each induce a clique; set \(q_2(G)=0\) if no such partition exists. The
same preprint~\cite[Theorem 5.1]{jiang-complete-first-factor} also treats
complete joins whose factors are all nonempty and noncomplete. Its Theorem 5.1 explicitly
does not classify joins containing both a complete and a noncomplete
factor. Thus the mixed join \(K_r+T\), with \(T\) a noncomplete tree,
lies outside that theorem; it is also different from the lexicographic
product \(K_m\circ T\).

One subfamily of the mixed-join problem is already known. The fan graph
is \(F_n=K_1+P_n\), and Tian, Dokyeesun, and Klav\v{z}ar~\cite[arXiv v2]{tian-et-al-removal}
proved, for \(n\ge 4\), that

\nopagebreak[4]
\[
\operatorname{gp}_d(F_n)
=\left\lfloor\frac{2(n+1)}{3}\right\rfloor
=\left\lceil\frac{2n}{3}\right\rceil.
\]

Consequently, the fan case is prior work and is used here only as a
consistency check. The formula above is the one displayed in arXiv v2
of~\cite{tian-et-al-removal}; no argument below relies on the different
rounding convention displayed in an earlier author-hosted version.

This paper studies \(K_r+T\) for \(r\ge 1\) and for trees \(T\) of order
at least three. For a tree \(T\), define

\[
\beta(T)=\max\bigl\{|X|:\ \Delta(T[X])\le 1
\text{ and } |N_T(x)\setminus X|\le 1\text{ for every }x\in X\bigr\}.
\]

The main structural argument separates dual general-position sets that
meet the universal clique \(V(K_r)\) from those that avoid it. The first
branch is controlled by a two-clique partition of the tree, while the
second is exactly the optimization problem defining \(\beta(T)\). This
yields

\[
\operatorname{gp}_d(K_r+T)=
\begin{cases}
r+2, & T\in\{P_3,P_4\},\\
\beta(T), & \text{otherwise}.
\end{cases}
\]

Sections~\ref{sec:structure}--\ref{sec:algorithm} give a self-contained proof of both branches, a local
characterization of the \(\beta\)-feasible sets, and a linear-time
dynamic program that reconstructs a maximum set. Section~\ref{sec:examples} gives worked
families, and Section~\ref{sec:computation} records two bounded checks with
distinct verification logic. The computations are supporting evidence rather
than proof. The precise comparison with the closest product, join, and fan
results is given above; no claim about graph classes outside the stated scope
is needed for the results proved below.

\section{Preliminaries}\label{sec:preliminaries}

All graphs in this paper are finite and simple. For a graph \(G\), its
vertex and edge sets are denoted by \(V(G)\) and \(E(G)\). If
\(X\subseteq V(G)\), then \(G[X]\) is the subgraph induced by \(X\), and
\(G-X\) abbreviates \(G[V(G)\setminus X]\). The open neighborhood of a
vertex \(x\) is \(N_G(x)\), and \(\Delta(G)\) denotes the maximum degree
of \(G\); by convention, the empty graph has maximum degree zero. A
complete graph of order \(r\) is denoted by \(K_r\), and
\(P_n\) denotes the path of order \(n\).

Let \(G\) be connected. The distance \(d_G(u,v)\) is the minimum length
of a \(u,v\)-path, where length means the number of edges. A path of
length \(d_G(u,v)\) is a \emph{\(u,v\)-geodesic}. There may be more than
one geodesic between the same two vertices. The interval between \(u\)
and \(v\) is

\[
I_G(u,v)=\{w\in V(G):w\text{ lies on some }u,v\text{-geodesic}\}.
\]

A set \(Y\subseteq V(G)\) is \emph{convex} if \(I_G(u,v)\subseteq Y\)
for every \(u,v\in Y\); equivalently, every geodesic between two
vertices of \(Y\) lies entirely in \(G[Y]\).

For \(X\subseteq V(G)\), a pair \(u,v\in V(G)\) is called
\emph{\(X\)-positionable} if

\[
I_G(u,v)\cap X\subseteq\{u,v\}.
\]

Thus no internal vertex of any \(u,v\)-geodesic belongs to \(X\). The
set \(X\) is a \emph{general-position set} if every pair of vertices in
\(X\) is \(X\)-positionable. Equivalently, no three distinct vertices of
\(X\) lie on a common geodesic.

A set \(X\subseteq V(G)\) is a \emph{dual general-position set} if every
pair whose two endpoints both lie in \(X\), or both lie in
\(V(G)\setminus X\), is \(X\)-positionable. The first family of pairs
says exactly that \(X\) is a general-position set. By Tian and Klav\v{z}ar
\cite[Theorem 3.1]{tian-klavzar-variety}, the second family is equivalently expressed by
convexity of the complement. We will therefore use the following
criterion throughout:

\nopagebreak[4]
\[
X\text{ is dual general position in }G
\quad\Longleftrightarrow\quad
X\text{ is general position and }G-X\text{ is convex}.
\]

\Needspace{6\baselineskip}
The \emph{dual general position number} is

\[
\operatorname{gp}_d(G)=
\max\{|X|:X\subseteq V(G)\text{ is a dual general-position set}\}.
\]

For vertex-disjoint graphs \(G\) and \(F\), their \emph{complete join}
\(G+F\) is obtained from their disjoint union by adding every edge with
one endpoint in \(V(G)\) and the other in \(V(F)\). In the rest of the
note, \(r\ge 1\), \(T\) is a tree of order at least three, and

\[
C=V(K_r),\qquad H=K_r+T.
\]

Every vertex of \(C\) is adjacent to every other vertex of \(H\). Two
vertices on the tree side are adjacent in \(H\) exactly when they are
adjacent in \(T\). Because such a tree is noncomplete, \(H\) has
diameter two.

We use two parameters on the tree side. First, define

\[
q_2(T)=\max\{|A|:V(T)=A\mathbin{\dot\cup}B,
\ T[A]\text{ and }T[B]\text{ are complete}\},
\]

and put \(q_2(T)=0\) if no such partition exists. The two classes may be
interchanged, so \(q_2(T)\) records the largest possible class in a
partition of the tree into two induced cliques. Second, define

\[
\beta(T)=\max\bigl\{|X|:X\subseteq V(T),\ \Delta(T[X])\le 1,
\ |N_T(x)\setminus X|\le 1\text{ for every }x\in X\bigr\}.
\]

The first constraint in the definition of \(\beta(T)\) says that a
selected vertex has at most one selected tree neighbor. The second says
that it also has at most one unselected tree neighbor. These are
conditions on the tree \(T\), not on the additional edges of the join
\(H\).

\section{The two-branch structural theorem}\label{sec:structure}

We first classify dual general-position sets that contain at least one
vertex of the universal clique \(C\). This is the \emph{apex-meeting
branch}.

\subsection{Sets meeting the universal clique}\label{sec:apex-meeting}

\begin{lemma}\label{lem:apex-meeting}
Let \(X\subseteq V(H)\) satisfy
\(X\cap C\ne\varnothing\), and put

\[
S=X\cap V(T).
\]

Then \(X\) is a dual general-position set of \(H\) if and only if both
\(T[S]\) and \(T[V(T)\setminus S]\) are complete.
\end{lemma}

\begin{proof}
Choose a vertex \(c\in X\cap C\). Suppose first that
\(X\) is a dual general-position set. If two vertices \(u,v\in S\) were
nonadjacent in \(T\), then they would also be nonadjacent in \(H\).
Since \(c\) is adjacent to every vertex of \(H\), the path \(u,c,v\)
would be a \(u,v\)-geodesic of length two containing three vertices of
\(X\). This would contradict the general-position property of \(X\).
Hence \(T[S]\) is complete.

Now suppose that two vertices \(u,v\in V(T)\setminus S\) were
nonadjacent in \(T\). Again \(u,c,v\) would be a \(u,v\)-geodesic of
length two. Its endpoints would lie in \(H-X\), whereas its internal
vertex \(c\) would lie in \(X\). Consequently \(H-X\) would not be
convex, a contradiction. Thus \(T[V(T)\setminus S]\) is complete.

Conversely, suppose that both induced tree subgraphs are complete. The
set \(X\) is then a clique: its vertices in \(C\) are pairwise adjacent,
its tree part \(S\) induces a clique, and every vertex of \(C\) is
adjacent to every vertex of \(T\). Therefore \(X\) is in general
position. The complement \(H-X\) is also a clique, because
\(C\setminus X\) is a clique, \(T[V(T)\setminus S]\) is a clique, and
all edges between these two parts are present. Every clique is convex,
so the criterion from Section~\ref{sec:preliminaries} shows that \(X\) is a dual
general-position set. This proves the equivalence.
\end{proof}

The converse also covers the allowed boundary case \(X\cap C=C\),
equivalently \(C\setminus X=\varnothing\). In the
present scope, \(T\) is noncomplete, so the two tree parts in Lemma~\ref{lem:apex-meeting}
cannot be empty in a feasible apex-meeting set: otherwise the other tree
part would induce all of \(T\).

\begin{corollary}\label{cor:apex-meeting-maximum}
A dual general-position set meeting \(C\) exists
if and only if \(q_2(T)>0\). When this branch exists, its maximum
cardinality is

\[
r+q_2(T).
\]
\end{corollary}

\begin{proof}
By Lemma~\ref{lem:apex-meeting}, the tree part \(S\) of such a set and its
complement form a partition of \(V(T)\) into two induced cliques. Thus
the branch can exist only if \(q_2(T)>0\). Conversely, if
\(V(T)=A\mathbin{\dot\cup}B\) is a partition into two induced cliques,
then Lemma~\ref{lem:apex-meeting} shows that \(C\cup A\) is a dual general-position set
meeting \(C\). This proves the existence equivalence.

For every set in this branch, \(|X\cap C|\le r\) and \(|S|\le q_2(T)\),
so \(|X|\le r+q_2(T)\). Choose a two-clique partition whose larger
designated class has size \(q_2(T)\) and include that class together
with all \(r\) vertices of \(C\). Lemma~\ref{lem:apex-meeting} gives a dual
general-position set of size \(r+q_2(T)\), proving that the bound is
attained.
\end{proof}

\subsection{Sets avoiding the universal clique}\label{sec:apex-avoiding}

We next consider sets contained entirely in the tree side. In this
branch the whole universal clique belongs to the complement.

\begin{lemma}\label{lem:apex-avoiding}
Let \(X\subseteq V(T)\), equivalently
\(X\cap C=\varnothing\). Then \(X\) is a dual general-position set of
\(H\) if and only if

\[
\Delta(T[X])\le 1
\]

and

\[
|N_T(x)\setminus X|\le 1\qquad\text{for every }x\in X.
\]
\end{lemma}

\begin{proof}
Recall that \(H\) has diameter two. We first examine the
general-position condition. Suppose that some \(x\in X\) has two
distinct neighbors \(u,v\in X\) in the tree. Since a tree contains no
triangle, \(u\) and \(v\) are nonadjacent. Hence \(u,x,v\) is a
\(u,v\)-geodesic of length two in \(H\), and all three of its vertices
belong to \(X\). Thus \(X\) is not in general position.

Conversely, suppose that \(X\) is not in general position. Then some
geodesic contains three distinct vertices of \(X\). Because \(H\) has
diameter two, these three vertices must occur as \(u,x,v\) on a
length-two \(u,v\)-geodesic, with \(x\) as its internal vertex. All
three lie in \(V(T)\), and adjacency between two tree-side vertices in
\(H\) is the same as adjacency in \(T\). Therefore \(u\) and \(v\) are
two selected tree neighbors of \(x\). We have proved that \(X\) is in
general position if and only if no selected vertex has two selected tree
neighbors, which is exactly the condition \(\Delta(T[X])\le 1\).

It remains to characterize convexity of the complement. Put
\(Y=V(H)\setminus X\). Since \(X\subseteq V(T)\), the entire clique
\(C\) lies in \(Y\). Any pair of vertices of \(Y\) with at least one
endpoint in \(C\) is adjacent. Likewise, two tree-side vertices of \(Y\)
that are adjacent in \(T\) are adjacent in \(H\). Such adjacent pairs
have geodesics of length one and cannot witness a failure of convexity.
Consequently \(Y\) can fail to be convex only when two nonadjacent
unselected tree vertices have a geodesic of length two whose internal
vertex lies in \(X\).

Because every vertex of \(X\) is on the tree side, such a geodesic
exists exactly when some \(x\in X\) has two distinct neighbors
\(u,v\in V(T)\setminus X\). Indeed, if such \(u\) and \(v\) exist, then
they are nonadjacent because \(T\) has no triangle, so \(u,x,v\) is a
length-two geodesic with endpoints in \(Y\) and internal vertex outside
\(Y\). Conversely, every length-two geodesic witnessing nonconvexity has
a selected internal vertex adjacent in \(T\) to its two unselected
endpoints. It follows that \(Y\) is convex if and only if
\(|N_T(x)\setminus X|\le 1\) for every \(x\in X\).

Combining the general-position and convexity equivalences with the
criterion from Section~\ref{sec:preliminaries} proves the lemma.
\end{proof}

\begin{corollary}\label{cor:apex-avoiding-maximum}
The maximum cardinality of a dual
general-position set that avoids \(C\) is \(\beta(T)\).
\end{corollary}

\begin{proof}
Lemma~\ref{lem:apex-avoiding} says that the dual general-position sets in
this branch are exactly the subsets of \(V(T)\) satisfying the two
constraints in the definition of \(\beta(T)\). Taking the maximum of
their cardinalities gives \(\beta(T)\).
\end{proof}

\subsection{Combining the two branches}\label{sec:combining}

\begin{theorem}\label{thm:two-branch}
For every \(r\ge 1\) and every tree \(T\) of order
at least three,

\[
\operatorname{gp}_d(K_r+T)=
\begin{cases}
\beta(T),&q_2(T)=0,\\
\max\{\beta(T),r+q_2(T)\},&q_2(T)>0.
\end{cases}
\]
\end{theorem}

\begin{proof}
Every vertex set \(X\subseteq V(H)\) either meets \(C\)
or avoids \(C\), and the two cases are disjoint. If \(q_2(T)=0\),
Corollary~\ref{cor:apex-meeting-maximum} says that the apex-meeting branch is empty, while
Corollary~\ref{cor:apex-avoiding-maximum} gives an attainable maximum of \(\beta(T)\) in the
apex-avoiding branch. If \(q_2(T)>0\), both branches exist: their
attainable maxima are respectively \(r+q_2(T)\) by Corollary~\ref{cor:apex-meeting-maximum} and
\(\beta(T)\) by Corollary~\ref{cor:apex-avoiding-maximum}. Taking the larger of these two values
proves the formula.
\end{proof}

For trees, the occurrence of the second line of Theorem~\ref{thm:two-branch} has a
particularly short classification.

\begin{lemma}\label{lem:q2-trees}
If \(T\) is a tree of order at least three, then

\[
q_2(T)=
\begin{cases}
2,&T\in\{P_3,P_4\},\\
0,&\text{otherwise}.
\end{cases}
\]
\end{lemma}

\begin{proof}
A tree contains no triangle. Hence every clique in \(T\)
has at most two vertices. If \(V(T)\) can be partitioned into two sets
that each induce a clique, then \(|V(T)|\le 4\).

There is only one tree of order three up to isomorphism, namely \(P_3\).
Its two adjacent vertices and its remaining vertex form a two-clique
partition, so \(q_2(P_3)=2\).

At order four, both classes in a two-clique partition must have size two
and therefore must be edges. Thus the two classes form a perfect
matching. The two trees of order four are \(P_4\) and \(K_{1,3}\). The
path \(P_4\) has a perfect matching, obtained by taking its first and
last edges, whereas two disjoint edges cannot occur in \(K_{1,3}\)
because every edge contains its center. Consequently \(q_2(P_4)=2\) and
\(q_2(K_{1,3})=0\). Trees of larger order cannot admit such a partition,
proving the lemma.
\end{proof}

\Needspace{6\baselineskip}
We also need the two corresponding values of \(\beta\).

\begin{lemma}\label{lem:beta-small-paths}
We have \(\beta(P_3)=2\) and \(\beta(P_4)=3\).
\end{lemma}

\begin{proof}
Let the vertices of \(P_3\) occur in the order
\(v_1,v_2,v_3\). The set \(\{v_1,v_3\}\) is \(\beta\)-feasible: it is
independent, and each of its vertices has only the neighbor \(v_2\)
outside the set. Hence \(\beta(P_3)\ge 2\). The full vertex set is not
feasible because its induced subgraph has maximum degree two. Since it
is the only set of size three, \(\beta(P_3)=2\).

Now let \(v_1,v_2,v_3,v_4\) be the vertices of \(P_4\) in path order.
The set \(\{v_1,v_2,v_4\}\) is \(\beta\)-feasible: its induced graph
consists of the edge \(v_1v_2\) and the isolated vertex \(v_4\), while
\(v_2\) and \(v_4\) each have exactly one neighbor outside the set and
\(v_1\) has none. Thus \(\beta(P_4)\ge 3\). The full vertex set is not
feasible because its induced subgraph has maximum degree two, so
\(\beta(P_4)\le 3\). Therefore \(\beta(P_4)=3\).
\end{proof}

Combining Theorem~\ref{thm:two-branch} with Lemmas~\ref{lem:q2-trees}
and~\ref{lem:beta-small-paths} yields the announced form.

\begin{corollary}\label{cor:main-formula}
For every \(r\ge 1\) and every tree \(T\) of
order at least three,

\[
\operatorname{gp}_d(K_r+T)=
\begin{cases}
r+2,&T\in\{P_3,P_4\},\\
\beta(T),&\text{otherwise}.
\end{cases}
\]
\end{corollary}

\begin{proof}
If \(T\notin\{P_3,P_4\}\), Lemma~\ref{lem:q2-trees} gives \(q_2(T)=0\),
so the first line of Theorem~\ref{thm:two-branch} yields
\(\operatorname{gp}_d(K_r+T)=\beta(T)\). For \(T\in\{P_3,P_4\}\), the
second line gives

\[
\operatorname{gp}_d(K_r+T)=\max\{\beta(T),r+2\}.
\]

Here \(\beta(P_3)=2\) and \(\beta(P_4)=3\) by Lemma~\ref{lem:beta-small-paths}, while
\(r+2\ge 3\). Thus the displayed maximum is \(r+2\) in both cases.
\end{proof}

\section{The tree parameter \texorpdfstring{$\beta(T)$}{beta(T)}}\label{sec:beta}

Call a set \(X\subseteq V(T)\) \emph{\(\beta\)-feasible} if it satisfies
the two constraints in the definition of \(\beta(T)\). Those constraints
admit the following equivalent description using only the degrees in the
original tree and the labels of the immediate neighbors.

\begin{proposition}\label{prop:local-characterization}
A set \(X\subseteq V(T)\) is
\(\beta\)-feasible if and only if both of the following conditions hold:

\begin{enumerate}
\def\labelenumi{\arabic{enumi}.}
\tightlist
\item
  no vertex of degree at least three in \(T\) belongs to \(X\); and
\item
  every vertex \(x\in X\) of degree two in \(T\) has exactly one
  neighbor in \(X\).
\end{enumerate}

Selected vertices of degree zero or one have no additional restriction.
In particular, the statement includes these boundary degrees even though
a degree-zero vertex does not occur under the standing assumption that
\(T\) has order at least three.
\end{proposition}

\begin{proof}
For each selected vertex \(x\in X\), put

\[
s_X(x)=|N_T(x)\cap X|,
\qquad
u_X(x)=|N_T(x)\setminus X|.
\]

Then

\[
s_X(x)+u_X(x)=\deg_T(x).
\]

The condition \(\Delta(T[X])\le 1\) says exactly that \(s_X(x)\le 1\)
for every \(x\in X\), while the second constraint in the definition of
\(\beta(T)\) says that \(u_X(x)\le 1\). Hence a selected vertex can have
degree at most two in \(T\). If its degree is two, the two nonnegative
integers \(s_X(x)\) and \(u_X(x)\) sum to two and are both at most one,
so they must both equal one. This proves the necessity of conditions 1
and 2.

Conversely, assume conditions 1 and 2. If \(x\in X\) has degree two,
then it has one selected and one unselected neighbor, so
\(s_X(x)=u_X(x)=1\). If \(x\in X\) has degree zero or one, then both
neighbor counts are automatically at most one. Condition 1 excludes
every other degree for a selected vertex. Thus \(s_X(x)\le 1\) and
\(u_X(x)\le 1\) for every \(x\in X\), which are precisely the two
defining constraints of a \(\beta\)-feasible set.
\end{proof}

This equivalence concerns every feasible set, not only the maximum ones.
One immediate consequence gives a useful general lower bound.

\begin{corollary}\label{cor:leaves}
If \(L(T)\) is the set of leaves of a tree
\(T\), then \(L(T)\) is \(\beta\)-feasible. Consequently,

\[
\beta(T)\ge |L(T)|.
\]
\end{corollary}

\begin{proof}
Every vertex in \(L(T)\) has degree one. Proposition~\ref{prop:local-characterization}
therefore imposes no additional condition on the set consisting of all
leaves. The inequality follows from the definition of \(\beta(T)\).
\end{proof}

The first constraint \(\Delta(T[X])\le 1\) alone defines a weaker
maximum-degree-one problem. Every \(\beta\)-feasible set satisfies that
constraint, but the converse is false. In \(K_{1,3}\), for example, the
center together with one leaf induces a single edge and satisfies the
weaker constraint. It is not \(\beta\)-feasible, because the selected
center has two unselected neighbors. Thus a result or algorithm for the
weaker problem cannot simply be substituted for \(\beta(T)\) without
enforcing the additional outside-neighbor constraint. The
bounded literature audit did not identify an established name for this
stricter parameter; this paper makes no claim that the parameter itself
is new.

\section{Linear-time computation and reconstruction}\label{sec:algorithm}

We now give a dynamic program that computes \(\beta(T)\) and
reconstructs an attaining set. Root \(T\) at an arbitrary vertex
\(\rho\). For a vertex \(v\), let \(T_v\) be the subtree induced by
\(v\) and all its descendants, and let \(\operatorname{Ch}(v)\) be the
set of children of \(v\). We use the label \(1\) for a selected vertex
and \(0\) for an unselected vertex.

For a nonroot vertex \(v\) and labels \(a,b\in\{0,1\}\), define
\(F_v(a,b)\) as the largest number of selected vertices in \(T_v\) among
all labelings with the following properties:

\begin{itemize}
\tightlist
\item
  \(v\) has label \(a\);
\item
  the parent of \(v\) has label \(b\); and
\item
  every selected vertex in \(T_v\) satisfies the two constraints
  defining a \(\beta\)-feasible set, where the specified parent label is
  used when checking the constraint at \(v\).
\end{itemize}

For the root, write \(F_\rho(a,\bot)\), where \(\bot\) records that no
parent exists. An impossible state has value \(-\infty\).

Fix a state at \(v\) and assign a proposed label \(t_u\in\{0,1\}\) to
every child \(u\in\operatorname{Ch}(v)\). Put

\[
s=\sum_{u\in\operatorname{Ch}(v)}t_u,
\qquad
d=|\operatorname{Ch}(v)|+\mathbf 1_{\{v\ne\rho\}},
\]

and define

\[
p=
\begin{cases}
0,&v=\rho,\\
b,&v\ne\rho.
\end{cases}
\]

Thus \(d=\deg_T(v)\), while \(p+s\) is the number of selected neighbors
of \(v\). If \(a=1\), the proposed child labels are admissible exactly
when

\nopagebreak[4]
\begin{equation}\label{eq:local-state-constraint}
p+s\le 1
\qquad\text{and}\qquad
d-(p+s)\le 1.
\end{equation}

The first inequality bounds the selected neighbors of \(v\) and the
second bounds its unselected neighbors. If \(a=0\), there is no
constraint at \(v\) itself, because the definition of \(\beta\) imposes
the two neighbor bounds only on selected vertices.

The recurrence is therefore

\begin{equation}\label{eq:dp-recurrence}
F_v(a,b)
=a+
\max_{(t_u:u\in\operatorname{Ch}(v))}
\sum_{u\in\operatorname{Ch}(v)}F_u(t_u,a),
\end{equation}

where all child-label vectors are allowed when \(a=0\), and only the
vectors satisfying \eqref{eq:local-state-constraint} are allowed when \(a=1\). For the root, the
same formula is used with \(b=\bot\) and \(p=0\). If the admissible
family is empty or a proposed vector uses an impossible child state, its
contribution is \(-\infty\). The answer is

\begin{equation}\label{eq:dp-answer}
\beta(T)=\max\{F_\rho(0,\bot),F_\rho(1,\bot)\}.
\end{equation}

\begin{theorem}\label{thm:dp-correctness}
The recurrence \eqref{eq:dp-recurrence} computes the exact values of
all states, and \eqref{eq:dp-answer} equals \(\beta(T)\). If one maximizing child-label
vector is stored for every finite state, a maximum \(\beta\)-feasible
set can be reconstructed.
\end{theorem}

\begin{proof}
We prove the state claim by induction on the height of
\(v\) in the rooted tree. If \(v\) is a leaf, it has no child subtrees.
When \(v\) is selected, \eqref{eq:local-state-constraint} checks exactly the label of its parent, if
one exists; when \(v\) is unselected, it has no constraint of its own.
Thus \eqref{eq:dp-recurrence} lists exactly the feasible labelings of the one-vertex
subtree and counts its selected vertex by the term \(a\).

Now assume that the claim holds for every child of \(v\). Fix the label
\(a\) of \(v\) and, if \(v\) is not the root, the parent label \(b\).
Once the child labels are fixed, different child subtrees have no edges
between them. By the induction hypothesis, \(F_u(t_u,a)\) is the exact
largest contribution from \(T_u\) under its two boundary labels.
Therefore the best total contribution from the child subtrees is the sum
appearing in \eqref{eq:dp-recurrence}.

It remains only to check the constraint at \(v\). If \(a=0\), no such
constraint exists, so every combination of feasible child states is
allowed. If \(a=1\), then \(p+s\) and \(d-(p+s)\) are exactly the
selected- and unselected-neighbor counts of \(v\). Hence \eqref{eq:local-state-constraint} accepts
precisely the combinations for which \(v\) satisfies the two \(\beta\)
constraints. The recurrence thus neither omits a feasible labeling nor
includes an infeasible one, and the additional term \(a\) counts \(v\)
exactly once. This proves the state claim by induction.

At the root, every labeling of the whole tree has exactly one of the two
root labels, so taking the maximum in \eqref{eq:dp-answer} gives precisely
\(\beta(T)\). To reconstruct an attaining set, choose a maximizing root
label. At each visited state, add \(v\) to the set if \(a=1\), read its
stored maximizing child-label vector, and continue to each child \(u\)
in state \((t_u,a)\). These boundary labels are compatible by
construction, and the preceding induction shows that every stored child
state attains its recorded value. The resulting set is therefore
\(\beta\)-feasible and has cardinality \(\beta(T)\).
\end{proof}

\begin{proposition}\label{prop:dp-complexity}
The value computation and the reconstruction
can both be implemented in \(O(|V(T)|)\) time using \(O(|V(T)|)\) space.
\end{proposition}

\begin{proof}
In a state with \(a=0\), the child labels are
independent, so for each child \(u\) we choose the larger of
\(F_u(0,0)\) and \(F_u(1,0)\). The work is therefore proportional to the
number of children. In a state with \(a=1\), \eqref{eq:local-state-constraint} is impossible when
\(d\ge 3\), because the two bounded neighbor counts sum to \(d\). Every
remaining selected state has at most two children, so at most four
child-label vectors need to be examined. There are only constantly many
states per vertex. Summing the child counts over all vertices gives
\(|V(T)|-1\), and hence the bottom-up computation takes linear time.

Storing a constant number of state values per vertex and one chosen
label for each relevant parent--child state uses linear space. The
top-down reconstruction visits every vertex once and therefore also
takes linear time.
\end{proof}

\section{Known subfamilies and worked examples}\label{sec:examples}

The examples below illustrate the local structure of \(\beta(T)\) and
the two different branches of the main theorem. Their values are derived
from the preceding proofs; no computation is being used as proof.

\subsection{Stars and the preserved counterexample}\label{sec:stars}

Let \(T=K_{1,k}\) with center \(z\) and \(k\ge 3\) leaves.
Proposition~\ref{prop:local-characterization} forces \(z\) to be
unselected because \(\deg_T(z)=k\ge 3\). On the
other hand, the set of all \(k\) leaves is \(\beta\)-feasible. It
follows that

\[
\beta(K_{1,k})=k.
\]

The star \(K_{1,k}\) is neither \(P_3\) nor \(P_4\) when \(k\ge 3\), so
Corollary~\ref{cor:main-formula} gives

\begin{equation}\label{eq:star-family}
\operatorname{gp}_d(K_r+K_{1,k})=k
\qquad (r\ge 1,\ k\ge 3).
\end{equation}

Thus the value is independent of the size of the universal clique. In
particular,

\[
\operatorname{gp}_d(K_r+K_{1,3})=3
\]

for every \(r\ge 1\). This is the counterexample retained from the
feasibility stage: \(q_2(K_{1,3})=0\), so the unsupported extrapolation
\(r+q_2(T)\) would give \(4\) when \(r=4\), whereas the correct value is
\(3\). The maximum set in the apex-avoiding branch consists of the three
leaves.

\subsection{Once-subdivided stars}\label{sec:subdivided-stars}

Let \(S_k\) be obtained from \(K_{1,k}\) by subdividing every edge
exactly once, where \(k\ge 3\). Write \(z\) for the center, \(u_i\) for
the degree-two vertex on the \(i\)th arm, and \(v_i\) for the
corresponding leaf. Again \(z\) is forced to be unselected. The set

\[
X=\{u_i,v_i:1\le i\le k\}=V(S_k)\setminus\{z\}
\]

is \(\beta\)-feasible: every selected \(u_i\) has the selected neighbor
\(v_i\) and the unselected neighbor \(z\), while each \(v_i\) has degree
one. Since the only remaining vertex is the forced-unselected center,
this set is maximum and

\[
\beta(S_k)=2k.
\]

The tree \(S_k\) has order \(2k+1\ge 7\), so \(q_2(S_k)=0\) by
Lemma~\ref{lem:q2-trees}. Hence

\begin{equation}\label{eq:subdivided-star-family}
\operatorname{gp}_d(K_r+S_k)=2k
\qquad (r\ge 1,\ k\ge 3).
\end{equation}

This family also has a value independent of \(r\), but unlike an
ordinary star its maximum set contains both leaves and degree-two
vertices.

\subsection{Paths and the known fan subfamily}\label{sec:paths}

We first determine \(\beta\) on paths directly from Proposition~\ref{prop:local-characterization}.

\begin{proposition}\label{prop:beta-path}
For every \(n\ge 3\),

\[
\beta(P_n)=\left\lceil\frac{2n}{3}\right\rceil.
\]
\end{proposition}

\begin{proof}
List the path vertices in order as \(v_1,\ldots,v_n\). A
\(\beta\)-feasible set cannot contain three consecutive vertices: the
middle one would be a selected degree-two vertex with two selected
neighbors, contrary to Proposition~\ref{prop:local-characterization}. Partitioning the path order into
consecutive blocks of three, followed by a remainder of size zero, one,
or two, therefore gives

\[
|X|\le
\begin{cases}
2q,&n=3q,\\
2q+1,&n=3q+1,\\
2q+2,&n=3q+2.
\end{cases}
\]

These three bounds equal \(\lceil 2n/3\rceil\). They are attained by the
binary selection words

\[
(110)^q,\qquad (110)^q1,\qquad (110)^q11,
\]

respectively, where a \(1\) means that the corresponding path vertex is
selected. In each word, every selected internal vertex has exactly one
selected neighbor, while the two endpoints have degree one.
Proposition~\ref{prop:local-characterization} therefore shows that the
indicated sets are \(\beta\)-feasible.
\end{proof}

Corollary~\ref{cor:main-formula} now yields, for every \(r\ge 1\) and \(n\ge 3\),

\begin{equation}\label{eq:path-family}
\operatorname{gp}_d(K_r+P_n)=
\begin{cases}
r+2,&n\in\{3,4\},\\
\left\lceil\dfrac{2n}{3}\right\rceil,&n\ge 5.
\end{cases}
\end{equation}

When \(r=1\), the graph \(K_1+P_n\) is the fan \(F_n\). For \(n=4\), the
first line of \eqref{eq:path-family} gives \(3=\lceil 8/3\rceil\), and for \(n\ge 5\) the
second line applies. Thus, for every \(n\ge 4\), the present theorem
recovers

\[
\operatorname{gp}_d(F_n)
=\left\lceil\frac{2n}{3}\right\rceil
=\left\lfloor\frac{2(n+1)}{3}\right\rfloor.
\]

The equality of the ceiling and floor expressions follows immediately by
considering \(n\) modulo three. This is exactly the formula stated in arXiv
v2 of~\cite{tian-et-al-removal}, so the fan calculation is a consistency
check and not a novelty claim. Formula \eqref{eq:path-family} additionally makes clear
that for paths of order at least five, the value stays unchanged when
the single fan apex is replaced by an arbitrary clique \(K_r\); the
exceptional paths \(P_3\) and \(P_4\) belong to the apex-meeting branch
and retain their \(r+2\) dependence.

\section{Computational verification and reproducibility}\label{sec:computation}

The proofs in Sections~\ref{sec:structure}--\ref{sec:examples} do not depend on computation. The
calculations reported here test the implementation and search for small
counterexamples; finite agreement is supporting evidence and is not a
proof of any theorem.

\subsection{Two verification routes}\label{sec:independent-checks}

The project keeps the following two main comparisons separate. They use
distinct checking logic but share the audit driver and NetworkX-generated
tree instances.

First, \path{src/mixed_join_tree.py} implements the rooted-tree
recurrence from Section~\ref{sec:algorithm} and reconstructs one maximizing set. It does
not import the shortest-path checker. The audit compares its value with
a separate exhaustive search over all subsets satisfying the two
defining local constraints of \(\beta(T)\). The comparison covers every
nonisomorphic tree of orders 3 through 12 generated by NetworkX, a total
of 985 trees. For each tree, the audit also checks that the
reconstructed set has the reported size and satisfies the two local
constraints. Finally, it recomputes the value with every possible root
to test root invariance.

Second, \path{src/dual_gp_independent.py} provides a
definition-first checker that does not use the mixed-join theorem, the
\(\beta\) characterization, or the tree DP. It constructs \(K_r+T\)
explicitly, computes all-pairs distances by fresh breadth-first
searches, enumerates vertex subsets, and tests general position together
with convexity of the complement. The audit compares the maximum found
by this checker with the formula implemented from Corollary~\ref{cor:main-formula}. It
covers all 46 nonisomorphic trees of orders 3 through 8 and
\(r\in\{1,2,3,4\}\), giving 184 comparisons.

\Needspace{8\baselineskip}
The archived report \path{results/mixed_join_dp_audit.json}
records:

\begin{longtable}[]{@{}L{0.25\textwidth}L{0.40\textwidth}rr@{}}
\toprule
Check & Scope & Number checked & Failures \\
\midrule
\endhead
DP value versus exhaustive subset search & all nonisomorphic trees,
orders 3--12 & 985 & 0 \\
root-invariance comparison & every possible root of the same 985 trees &
11,003 & 0 \\
reconstructed maximum set & one reconstruction for each of the 985 trees
& 985 & 0 \\
mixed-join formula versus shortest-path definition & all nonisomorphic
trees, orders 3--8, and \(r=1,2,3,4\) & 184 & 0 \\
\bottomrule
\end{longtable}

The earlier target-selection screen in
\path{results/extension_feasibility_audit.json} contains 92
overlapping small mixed-join comparisons and the explicitly retained
\(K_{1,3}\) data. Those historical checks are not added to the table
totals, because doing so would double-count instances from the larger
audit.

\subsection{Environment and commands}\label{sec:environment}

The archived matrix was produced under CPython 3.13.5 and NetworkX
3.6.1. A fresh rerun on 29 August 2026 used Windows 11 build 26100,
CPython 3.13.5, NetworkX 3.6.1, and pytest 9.1.1. From the repository
root, the relevant PowerShell commands are

\begin{verbatim}
.\.venv\Scripts\python.exe -m pip install -r .\requirements.txt
.\.venv\Scripts\python.exe -m pytest -q .\tests
.\.venv\Scripts\python.exe .\experiments\audit_mixed_join_dp.py `
  --output .\results\mixed_join_dp_audit.json
\end{verbatim}

The fresh test run returned \texttt{29\ passed}. In addition to the
bounded matrices, the tests cover named values, invalid inputs,
reconstruction feasibility, and a path of order 2,500 to ensure that
reconstruction does not rely on Python recursion. The fresh audit rerun
reproduced all four zero-failure rows in the table above.

\Needspace{7\baselineskip}
For exact artifact identification, the SHA-256 hashes recorded on 29
August 2026 are:

\begin{longtable}[]{@{}L{0.42\textwidth}L{0.53\textwidth}@{}}
\toprule
Artifact & SHA-256 \\
\midrule
\endhead
\path{src/mixed_join_tree.py} &
\texttt{D3DDAB46\allowbreak 510B217B\allowbreak F9C442CC\allowbreak
57302953\allowbreak C4D0C29F\allowbreak 208DD7AD\allowbreak
E06381F2\allowbreak D90A8B9D} \\
\path{src/dual_gp_independent.py} &
\texttt{D7BA1C52\allowbreak 0DC6991E\allowbreak 78CB8BAE\allowbreak
609A598E\allowbreak 6BBF950D\allowbreak A03C1EAC\allowbreak
9BA4D30D\allowbreak 1EFFF231} \\
\path{experiments/}\newline\path{audit_extension_candidates.py} &
\texttt{A463CE20\allowbreak 2995000F\allowbreak 324BDB7F\allowbreak
90D2821B\allowbreak 1B97AB38\allowbreak 911E3E89\allowbreak
41DFFC38\allowbreak 1F8393F6} \\
\path{experiments/}\newline\path{audit_mixed_join_dp.py} &
\texttt{93BA42DB\allowbreak 963CAB3F\allowbreak 2EE34DAC\allowbreak
9CF36511\allowbreak FF6332A8\allowbreak 4543F44A\allowbreak
22ACD48B\allowbreak BB974D44} \\
\path{results/}\newline\path{mixed_join_dp_audit.json} &
\texttt{C4489D2A\allowbreak 7202AD14\allowbreak 13EED2FB\allowbreak
C551E6CB\allowbreak DFEB6FC5\allowbreak BB10FF32\allowbreak
F4899B74\allowbreak 7C3F4E90} \\
\bottomrule
\end{longtable}

The \path{requirements.txt} file lists the required packages but does
not pin exact versions. The explicit versions and hashes above are
therefore part of the reproducibility record. None of these finite
checks establishes correctness beyond the stated test ranges.

\section{Conclusion and limitations}\label{sec:conclusion}

For \(r\ge 1\) and a tree \(T\) of order at least three, dual
general-position sets in \(K_r+T\) split into two exhaustive classes. A
set meeting the universal clique exists precisely when the tree can be
partitioned into two induced cliques, and the maximum in that branch is
\(r+q_2(T)\). A set avoiding the universal clique is dual general
position precisely when it satisfies the two local constraints defining
\(\beta(T)\), and the maximum in that branch is \(\beta(T)\). Since a
tree in the present scope has a two-clique partition only when it is
\(P_3\) or \(P_4\), the combined formula is

\[
\operatorname{gp}_d(K_r+T)=
\begin{cases}
r+2,&T\in\{P_3,P_4\},\\
\beta(T),&\text{otherwise}.
\end{cases}
\]

The local characterization of \(\beta\) excludes selected vertices of
degree at least three and requires every selected degree-two vertex to
have exactly one selected neighbor. The rooted-tree dynamic program
turns this characterization into a linear-time, linear-space algorithm
that also reconstructs a maximum set. Stars, once-subdivided stars, and
paths illustrate the resulting structure. For \(r=1\), the path
calculation recovers the formula stated in arXiv v2 of
\cite{tian-et-al-removal}; that subfamily is prior work and is not presented as
a new result.

Several limitations are essential. Jiang's closest general result
\cite{jiang-complete-first-factor} was available as preprint v1.0.1 at the
literature cutoff of 29 August 2026. The bounded computations in
Section~\ref{sec:computation} support the implementation but do not prove the
theorem. The comparison with the known fan formula uses the display in arXiv
v2 of~\cite{tian-et-al-removal}, which has the rounding convention consistent
with formula~\eqref{eq:path-family}.

Finally, this paper treats only \(K_r+T\) with \(r\ge 1\) and
\(|V(T)|\ge 3\). It does not classify joins with an arbitrary first
factor, general mixed complete joins with several factors, or
lexicographic products \(P_n\circ T\). Those directions are outside the
present scope, and no claim about their resolution or novelty is made
here.

\section*{Statements and Declarations}

\noindent\textbf{Funding.}
The author received no funding for this work.

\medskip
\noindent\textbf{Competing interests.}
The author declares no competing interests.

\medskip
\noindent\textbf{Author contributions.}
Yi Yuteng is the sole author and assumes responsibility for the study, the
reported mathematical arguments and computations, and the final manuscript.

\medskip
\noindent\textbf{Data and code availability.}
The code and machine-readable audit outputs supporting the computational
checks are available from the corresponding author on reasonable request.

\medskip
\noindent\textbf{Acknowledgements.}
The author acknowledges the assistance of OpenAI Codex in the roles described
below.

\medskip
\noindent\textbf{Use of AI-assisted tools.}
OpenAI Codex was used as an AI-assisted tool for literature organization,
proof drafting, code development, computational checking, manuscript drafting,
and editorial revision. Codex is not an author. The author assumes full
responsibility for the final manuscript, including verification of the
mathematical claims, citations, code, and text.

\clearpage
\begin{thebibliography}{9}

\bibitem{tian-klavzar-variety}
J. Tian and S. Klav\v{z}ar, ``Variety of general position problems
in graphs,'' \emph{Bulletin of the Malaysian Mathematical Sciences
Society} 48 (2025), Article 5,
\href{https://doi.org/10.1007/s40840-024-01788-z}{doi:10.1007/s40840-024-\mbox{01788-z}}.

\bibitem{chandran-et-al-survey}
S. V. Ullas Chandran, S. Klav\v{z}ar, and J. Tuite, ``The General
Position Problem: A Survey,'' arXiv:2501.19385v5 (16 August 2026),
\href{https://doi.org/10.48550/arXiv.2501.19385}{doi:10.48550/arXiv.2501.19385}.

\bibitem{ghorbani-et-al-operations}
M. Ghorbani, H. R. Maimani, M. Momeni, F. R. Mahid, S. Klav\v{z}ar,
and G. Rus, ``The general position problem on Kneser graphs and on some
graph operations,'' \emph{Discussiones Mathematicae Graph Theory} 41
(2021), 1199--1213,
\href{https://doi.org/10.7151/dmgt.2269}{doi:10.7151/dmgt.2269}.

\bibitem{dokyeesun-et-al-products}
P. Dokyeesun, S. Klav\v{z}ar, D. Kuziak, and J. Tian, ``General
position problems in strong and lexicographic products of graphs,''
\emph{Computational and Applied Mathematics} 45 (2026), Article 97,
\href{https://doi.org/10.1007/s40314-025-03547-7}{doi:10.1007/s40314-025-03547-7}.

\bibitem{jiang-complete-first-factor}
W. Jiang, ``Dual General Position in Lexicographic Products with
a Complete First Factor,'' Zenodo preprint v1.0.1 (27 August 2026),
\href{https://doi.org/10.5281/zenodo.22116770}{doi:10.5281/zenodo.22116770}.

\bibitem{tian-et-al-removal}
J. Tian, P. Dokyeesun, and S. Klav\v{z}ar, ``On the variety of
general position problems under vertex and edge removal,''
\emph{Discrete Applied Mathematics} 388 (2026), 56--64,
\href{https://doi.org/10.1016/j.dam.2026.02.044}{doi:10.1016/j.dam.2026.02.044};
\href{https://arxiv.org/abs/2510.01294v2}{arXiv:2510.01294v2}.

\end{thebibliography}

\end{document}
